\section{Noções de Medida}

A moderna Teoria das Probabilidades é fundamenta utilizando-se conceitos da Teoria
da Medida, um ramo da Análise que generaliza as noções de comprimento, área e volume
para espaços abstratos associando-os a uma ``medida'', a partir da qual também
estende-se uma Teoria de Integração para esses espaços. Probabilidades são vistas,
portanto, como medidas das chances de ocorrência de um subconjunto de eventos
dentro de um universo de resultados possíveis.

Nesta seção, será realizada uma breve exposição dos principais enunciados sobre medida
a fim de somente fundamentar o desenvolvimento posterior dos conceitos associados à
Teoria de Probabilidades, que, por sua vez, garante a fundamentação da teoria e
prática estatística.

\subsection{Adendos de Análise}

O conceito de continuidade, estudado dentro da Análise, está atrelado à ideia de
poder aproximar uma função real a um determinado valor variando suavemente seu
argumento.

\begin{def:pontualmente continua}
Um função $f: \real \to \real$ é uniformemente contínua em $a \in \real$ se
\[
	\forall \epsilon > 0\, \exists \delta \in \real\, \forall x \in \real\,,
	| x - a | < \delta \rightarrow | f(x) - f(a) | < \epsilon
\]
\end{def:pontualmente continua}

Outra forma equivalente de definição de continuidade é
\[
	\lim_{x \to a} f(x) = \lim_{x \to a^+} f(x) = \lim_{x \to a^-} f(x) = f(a).
\]
Funções que não são contínuas em determinado ponto signfica que, a depender de onde nos
aproximamos, o valor para o qual a função converge se diferencia.

\begin{exp:pontualmente contínua}
A função $f(x) = x^2$ é contínua em todo ponto de seu domínio, já que
\[
	\lim_{x \to a} x^2 = a^2
\]
\end{exp:pontualmente contínua}

\begin{exp:pontualmente contínua neg}
Funções escadas são exemplo de funções não contínuas.
\end{exp:pontualmente contínua neg}

A próxima definição é uma versão mais forte da anterior.
A continuidade uniforme signigica que, independente da escolha de limite de erro $\epsilon$
para o valor da função, há um único $\delta$ tal que em qualquer lugar observa-se que,
se $|x - a| < \delta$, então $|f(x) - f(a)| < \epsilon$.

\begin{def:uniformemente continua}
Uma função $f: \real \to \real$ é uniformemente contínua se
\[
	\forall \epsilon > 0\, \exists \delta \in \real\, \forall x,a \in \real\,,
	| x - a | < \delta \rightarrow | f(x) - f(a) | < \epsilon
\]
\end{def:uniformemente continua}

Uma função que não é uniformemente contínua pode ter variações muita abruptas em alguns
intervalos, e mais suaves em outros.

\begin{exp:uniformemente contínua}
As funções $\sin x$, $\cos x$ e funções lineares são exemplos de absolutamente contínuas,
pois suas taxas de variação não se tornam infinitamente grandes conforme a função se
aproxima de qualquer ponto.
\end{exp:uniformemente contínua}

\begin{exp:uniformemente contínua neg}
A função $f(x) = 1/x$ não é uniformemente contínua, pois conforme se aproxima de $0$,
a sua variação se torna infinitamente grande.
\end{exp:uniformemente contínua neg}

\begin{thm:heine}
Seja $f$ uma função pontualmente contínua em $X$. Se $X$ é
compacto (fechado e limitado), então $f$ é uniformemente contínua em $X$.
\end{thm:heine}

\subsection{Introdução à Medida}

Antes de podermos medir determinados conjuntos, precisamos entender que tipos de
conjuntos podem ser medidos.
Com esse propósito, introduz-se o conceito de $\sigma$-álgebras, que define uma álgebra
sobre uma coleção de partes de um conjunto cujas operações definem conjuntos
\textbf{mensuráveis}.

\begin{def:sigma algebra}
Seja $\Omega$ um conjunto e $\mathcal{B} \subset \Parts{\Omega}$.
O conjunto $\mathcal{B}$ é uma $\sigma$-álgebra de $\Omega$ se, e somente se,
\begin{enumerate}
	\item[i] $\varnothing \in \mathcal{B}$.
	\item[ii] $\forall A \in \mathcal{B}, A^c \in \mathcal{B}$.
	\item[iii] $\forall A \in \mathcal{B},\ \bigcup_{i=1}^\infty A_i \in \mathcal{B}$.
\end{enumerate}
Se $\Omega$ é um espaço amostral, então $A \in \mathcal{B}$ é denominado evento de
$\Omega$.
\end{def:sigma algebra}

\begin{def:medida nula}
Um conjunto é dito de medida nula se para todo $\epsilon > 0$ existe uma cobertura
enumerável $\{I_k\}_1^n$ tal que a soma dos comprimentos de cada intervalo $I_k$ é
menor que $\epsilon$.
\end{def:medida nula}

\begin{def:mensuravel}
\label{def:mensuravel}
O par $(\Omega, \mathcal{B})$ é dito mensurável se $\mathcal{B}$ for uma $\sigma$-álgebra
de $\Omega$.
A função \(\mu: \mathcal{B} \to [0, + \infty)\) é uma medida para $(\Omega, \mathcal{B})$
	se satisfaz
	\begin{equation*}
		\mu(\varnothing) = 0 \quad \text{e} \quad \mu\left(\bigsqcup_{i = 1}^\infty A_i\right) = \sum_{i=1}^\infty \mu(A_i)
	\end{equation*}
	A tripla $(\Omega, \mathcal{B}, \mu)$ é chamado de espaço de medida.
	\end{def:mensuravel}

	\begin{def:medida de lesbergue}
	A medida $\lambda$ de Lesbergue para $(\real, \borel)$ satisfaz
	\[
	\lambda((a, b]) = b - a
\]
\end{def:medida de lesbergue}

\begin{thm:calculo medida}
Se $(\Omega, \mathcal{B}, \mu)$ é uma espaço de medida e $\{A_n\}$ é uma família
enumerável de elementos de $\mathcal{B}$, então
\begin{enumerate}
	\item[i.] Se $A_1 \subset A_2 \subset \ldots$, então
	      \[
		      \mu\left(\bigcup_{n=1}^\infty A_n\right) = \lim_{n \to \infty} \mu(A_n)
	      \]
	\item[ii.] Se $A_1 \supset A_2 \supset \ldots$, então
	      \[
		      \mu\left(\bigcap_{n=1}^\infty A_n\right) = \lim_{n \to \infty} \mu(A_n)
	      \]
	\item[iii.]
	      \[
		      \mu\left(\bigcup_{n=1}^\infty A_n\right) \le \sum_{n=1}^\infty \mu(A_n)
	      \]
\end{enumerate}
\end{thm:calculo medida}

\begin{thm:medida de subconjunto}
\label{thm:medida de subconjunto}
Seja $(\Omega, \mathcal{B}, \mu)$ espaço de medida, e $A$ e $B$ elemetos de $\mathcal{B}$.
\[
	A \subset B \rightarrow \mu(A) \le \mu(B)
\]
\end{thm:medida de subconjunto}

\begin{def:medida completa}
Uma medida $\mu$ para $(\Omega, \mathcal{B})$ é dita completa se
\[
	\forall S \in \mathcal{B}\, \forall A\,, \mu(S) = 0 \land (A \subset S \rightarrow A \in \mathcal{B})
\]
A medida $\mu$ é dita $\sigma$-finita se
\[
	\Omega = \bigcup_{n=1}^\infty A_n
\]
com $\mu(A_n) < \infty$ para todo $n \in \nat$.
\end{def:medida completa}

\begin{thm:lesbergue sigma finita}
A medida de Lesbergue é $\sigma$-finita.
\end{thm:lesbergue sigma finita}

\begin{def:funcao mensuravel}
Sejam $(\Omega, \mathbb{B})$ e $(\Omega', \mathbb{B}')$	espaços de medida. Uma função
$f: \Omega \to \Omega'$ é mensurável -- ou $(\mathcal{B}, \mathcal{B}')$-mensurável --
se
\[
	\forall B \in \mathcal{B}', f^{-1}(B) \in \mathcal{B}
\]
onde $f^{-1}$ denota a pré-imagem de $f$.
\end{def:funcao mensuravel}

\begin{thm:funcoes mensuraveis}
São verdadeiras as seguintes afirmações:
\begin{enumerate}
	\item[i.] Toda função contínua é mensurável.
	\item[ii.] Se $f$ e $g$ são funções mensuráveis, então são mensuráveis
	      \[
		      f \circ g, \quad f \cdot g, \quad \frac{f}{g}, g \ne 0
	      \]
	\item[iii.] Se $\{f_n\}$ é uma sequência de funções mensuráveis, então são mensuráveis
	      \[
		      \lim \sup f_n, \quad \lim \inf f_n, \quad \sup f_n, \quad \inf f_n
	      \]
	\item[iv.] O limite pontual de uma função mensurável é mensurável.
\end{enumerate}
\end{thm:funcoes mensuraveis}

\subsection{Integral de Lesbergue}

\begin{def:integral de lesbergue}
Seja $f$ uma função. A integral de Lesbergue de $f$ é definida por
\[
	\int f\, d \lambda = \int_\real f(x) \cdot I_A(x)\, dx
\]
\end{def:integral de lesbergue}

\begin{thm:integral de lesbergue}
A integral de Lesbergue satisfaz
\begin{enumerate}
	\item[i.] Linearidade
	      \[
		      \int (a f + b g)\, d \lambda = a \int f\, d \lambda + b \int g\, d\lambda
	      \]
	\item[ii.] Monotonicidade
	      \[
		      f \le g \rightarrow \int f\, d\lambda \le \int g\, d\lambda
	      \]
	\item[iii.]
	      \[
		      f \ge 0 \rightarrow \int_A f\, d\lambda \ge 0
	      \]
	\item[iv.] Se $A \cap B = \varnothing$, então
	      \[
		      \int_{A \cup B} f\, d\lambda = \int_A f\, d\lambda + \int_B f\, d\lambda
	      \]
	\item[v.]
	      \[
		      \int_A f\, d \lambda  = \int f I_A\, d\lambda
	      \]
\end{enumerate}
\end{thm:integral de lesbergue}

\begin{def:quase sempre}
Uma proposição $\varphi(x)$ vale quase sempre (q.s.) num conjunto mensurável $A$ se
\[
	\lambda(\{x \in A \mid \neg \varphi(x)\} = 0
\]
\end{def:quase sempre}

\begin{thm:0 quase sempre}
Temos $f = 0$ q.s. em $A$ se, e somente se, $f \ge 0$ e $\int_A f\, d\lambda = 0$.
\end{thm:0 quase sempre}

\begin{thm:riemann}
Se $f$ é Riemann integrável, então
\[
	\int_{[a,b]} f \, d \lambda = \int_a^b f(x)\, dx
\]
\end{thm:riemann}
